\documentclass[../../../include/open-logic-section]{subfiles}

\begin{document}

\olfileid{sth}{choice}{justifications}
\olsection{د انتخاب په اړه دننۍ کتنې}

نو پراخه پوښتنه دا ده چې ايا ښه ترتيب،
يا انتخاب، يا د ټولو سټونو د اندازې
پرتله کېدنه — هر يو يې واخلئ، توپير
نه کوي — توجيه کېدلے شي؟

دا يې د \emph{دننۍ} توجيې يوه هڅه
ده۔ مخکې په \olref[z][story]{sec}
کښې مو د مراتبي جوړښت څو اصلونه
وړاندې کړي وو۔ يو يې بيا يادول
ارزي:
\begin{enumerate}
	\item[] \stagesacc. د هر پړاو $S$ دپاره، او د هغو
	هر ډول سټونو دپاره چې د $S$ له پړاو
	\emph{مخکې} جوړ شوي وو: په $S$ پړاو
	کښې داسې سټ جوړېږي چې غړي يې
	عين هماغه سټونه دي۔ په
	پړاو~$S$ کښې نور هېڅ نه جوړېږي۔
\end{enumerate}
په حقيقت کښې ګڼو ليکوالانو ويلي چې
د انتخاب اصل د (همدې ته ورته) اصل
په وسيله توجيه کېدلے شي۔ د دې لارې
لنډه تشريح وړاندې کوو۔

له يوې کوچنۍ ساده پايلې پيل کوو چې
د انتخاب يو \emph{بل} همارز اصل
وړاندې کوي:

\begin{thm}[په $\ZF$ کښې]\ollabel{choiceset}
انتخاب له لاندې اصل سره همارز دے۔
که د $A$ !!{element}s دوه په دوه بېل
او ناخالي وي، نو داسې $C$ شته چې
$C
\cap x$ د هر $x \in A$ دپاره
يو غړيز سټ وي۔ (داسې $C$ ته د
$A$ دپاره {د انتخاب سټ} وايو۔)
\end{thm}

د دې پايلې ثبوت نېغ دے، او لوستونکي
ته يې د تمرين په توګه پرېږدو۔

\begin{prob}
\olref[sth][choice][justifications]{choiceset}
ثابت کړئ۔ که ستونزه لرئ، ثبوت په
\cite[pp.~242--3]{Potter2004} کښې
موندلے شئ۔
\end{prob}

بنسټيزه خبره دا ده چې د $A$ دپاره
د انتخاب سټ د $A$ دپاره د انتخاب
تابعې د ارزښتونو سټ دے۔ نو د انتخاب
د توجيې دپاره يوازې د هغه همارزه
بڼه، د انتخاب سټونو د شتون له مخې،
توجيه کولے شو۔ اوس همدا هڅه کوو۔

فرض کړئ د $A$ !!{element}s دوه په
دوه بېل او ناخالي دي۔ د \stageshier{}
له مخې (وګورئ \olref[z][story]{sec})
$A$ په کوم پړاو~$S$ کښې جوړېږي۔
پام وکړئ چې د $\bigcup A$ ټول
!!{element}s د $S$ پړاو څخه
مخکې شته۔ اوس د \stagesacc{}
له مخې، د پړاو~$S$ څخه مخکې
جوړ شويو \emph{هر ډول} سټونو
دپاره داسې سټ جوړېږي چې غړي
يې عين هماغه سټونه دي۔ په بله
وينا، د مخکينيو شته سټونو هره
\emph{ممکنه} ټولګه به په~$S$
کښې شتون ولري۔ خو په يقين سره
داسې څيزونه ټاکل \emph{ممکن}
دي چې د~$A$ دپاره د انتخاب سټ
ترې جوړ شي؛ دا د $\bigcup A$
يو ډېر ځانګړے فرعي سټ دے۔ نو
د اړتيا له مخې داسې د انتخاب
سټ شته۔

دا د دننيو دلايلو پر بنسټ د
انتخاب د توجيې \emph{ډېره}
لنډه هڅه وه۔ خو د دې فکر د لا
څېړلو دپاره د پوټر
(\citeyear[\S14.8]{Potter2004})
ښه پراختيا ولولئ۔

\end{document}